\section{Experimental Setup}
\label{sec:experimental_setup}

\subsection{Datasets}

The suite contains six fixed evaluation snapshots over five QA families:
biomedical control, clinical QA, and open-domain multi-hop QA
(Table~\ref{tab:datasets}).
It deliberately mixes corpus contracts: per-question bundles, per-question
source documents, and shared corpora.
The primary graph object is a passage-provenance KG built from the retrieval
corpus itself, since the question is whether \system{} preserves dense-RAG
answer quality while exposing facts, passages, relations, and retrieval state.
PubMedQA and MuSiQue use $n=100$; RealMedQA uses its full $n=230$ evaluable
set; HotpotQA, HotpotQA FullWiki, and 2WikiMultiHopQA use fixed $n=250$
subsets.

\begin{table*}[t]
\centering
\small
\caption{Evaluation datasets, corpus contracts, KG scale, and primary roles.
$h$ is the KG traversal depth; \emph{corpus} denotes the retrieval contract;
$n$ is the evaluated subset size.
$|E|$ and $|R|$ count entity nodes and typed entity--entity relations in the
dataset-scoped KG, queried from the persistent Neo4j stores used in the
reported runs.
\label{tab:datasets}}
\setlength{\tabcolsep}{3pt}
\begin{tabular}{@{}lllclrrrp{2.7cm}@{}}
\toprule
Dataset & Domain & Answer type & $h$ & Corpus & $n$ & $|E|$ & $|R|$ & Primary role \\
\midrule
PubMedQA
  & Biomedical & Yes/No/Maybe & 2 & abstract & 100 & 2{,}354 & 3{,}023 & Binary control; strong dense baseline \\
RealMedQA
  & Clinical & Free-text & 2 & shared (143) & 230 & 537 & 359 & Clinical shared-corpus grounding \\
HotpotQA \citep{yang2018hotpotqa}
  & Wikipedia & Free-text & 2 & bundle & 250 & 13{,}546 & 12{,}355 & Open-domain bridge diagnostic \\
HotpotQA FullWiki \citep{yang2018hotpotqa}
  & Wikipedia & Free-text & 2 & shared subset & 250 & 13{,}498 & 12{,}012 & Shared-corpus bridge stress test \\
2WikiMultiHopQA \citep{ho2020wikimultihopqa}
  & Wikipedia & Free-text & 2 & bundle & 250 & 8{,}328 & 7{,}187 & Bridge and comparison analysis \\
MuSiQue \citep{trivedi2022musique}
  & Wikipedia & Free-text & 4 & bundle & 100 & 15{,}440 & 24{,}858 & Hard multi-hop stress test \\
\bottomrule
\end{tabular}
\end{table*}

\textsc{RealMedQA} is the main shared-corpus clinical setting and
\textsc{HotpotQA FullWiki} a controlled Wikipedia shared-corpus stress
snapshot; the rest are closed source-document or bundle settings.
This is a conservative testbed for KG-RAG: dense retrieval sees small,
benchmark-filtered candidate sets and provides a demanding accuracy baseline.
\textsc{PubMedQA} is retained as a control; the open-domain datasets
stress bridge preservation, shared-corpus evidence selection, comparison
questions, and deeper compositional chains.

Throughout, hop count refers to the reasoning chain implied by the
question, not the number of passages supplied by the benchmark.
Explicit decomposition metadata are used where available, otherwise
dataset-specific conventions.
Hop-stratified analyses are reported as diagnostics rather than the main
table.

\paragraph{Scope.}
The experiments are fixed-subset diagnostic runs, not leaderboard submissions
or web-scale throughput benchmarks.
They support the paper's coarse family-level claims; fine-grained ranking of
neighbouring AUROC or AUREC values would require larger multi-seed sweeps.

\subsection{Model and Sampling}

All answer-generation experiments use GPT-4o-mini
\citep{openai2024gpt4omini}.
For answer-state uncertainty, $N=5$ responses are sampled at temperature
$T=1.0$.
These are low-budget black-box diagnostics: the API provides no reliable token
likelihoods, so entropy-style scores use sampled responses rather than decoder
probabilities.

KG construction is largely deterministic: entity and relation extraction use
temperature $0.0$, and benchmark passages are ingested one at a time with
chunk-level provenance.
The studied regime is stabilised entity-first retrieval with explicit fallback.

\paragraph{What is and is not isolated.}
The comparison holds the corpus, embeddings, and generation model fixed, but
does not fully isolate graph structure from retrieval stabilisation, fallback
routing, prompt organisation, or provenance grouping
(Section~\ref{sec:limitations}).
The runs report mean pairwise chunk overlap across the $N=5$ calls
(Appendix Table~\ref{tab:retrieval_overlap}); an archived 2WikiMultiHopQA
trace also supports the stability split in
Table~\ref{tab:within_adaptive_stability}.
The headline runs do not retain per-question entity- or path-overlap statistics
for every dataset, but a dedicated graph-state diversity run on HotpotQA-FullWiki
logs the full per-sample seed-entity, path, subgraph, and chunk overlap and is
analysed in Appendix~\ref{sec:appendix_diversity}.

\subsection{Correctness Labels}
\label{sec:correctness}

Correctness labels are binary.
Label-space tasks are normalised to their answer contract
(\texttt{yes}/\texttt{no}/\texttt{maybe}).
Free-text and factoid tasks use a reference-based semantic judge given the
question, gold answer, aliases where available, and model response; by default
the judge is GPT-4o-mini
(Appendix~\ref{sec:appendix_reproducibility}).
This is a benchmark answer-normalisation task rather than open-ended clinical
diagnosis: the judge decides reference-grounded semantic equivalence, for which
exact match would undercount aliases, paraphrases, and short explanatory
answers.
Same-model judging is a known limitation
\citep{zheng2023judging,panickssery2024llm}, so the labels are practical
reference-grounded judgements rather than an oracle.
An independent re-judge with a different model family (Llama-3.3-70B) on the
saved answers of every free-text run, including both headline RealMedQA runs,
HotpotQA, HotpotQA FullWiki, MuSiQue, and both 2WikiMultiHopQA runs, agrees
with the original labels on $92$--$99\%$ of answers
($\kappa = 0.52$--$0.97$, the low end being the near-ceiling-accuracy
RealMedQA adaptive run where few wrong answers make $\kappa$ unstable) and
leaves every central contrast unchanged (Section~\ref{sec:limitations}).

Two conventions matter.
Provider-side failures are recorded explicitly, so results report both
\textbf{raw accuracy} and \textbf{clean accuracy}, the latter excluding
generation failures (Appendix Table~\ref{tab:generation_failures}).
For free-text datasets, prompts elicit short explanatory answers rather than
extractive spans; semantic correctness is therefore the headline answer metric,
with EM/F1 retained only in the artefacts.

\subsection{Evaluation Metrics}

The main uncertainty metric is \textbf{AUROC}, the ranking quality for
separating correct from incorrect answers.
\textbf{AUREC} is logged as a secondary selective-prediction audit metric
\citep{geifman2017selective}, but not expanded into a parallel results table.
Raw accuracy, clean accuracy, retrieval overlap, and per-metric compute time
are also logged; the marginal cost of each diagnostic, including SEU's
per-chunk NLI calls and GPS, is reported in
Appendix Table~\ref{tab:compute_cost}.
All main figures report point estimates from one fixed subset per dataset.
With mixed sample sizes ($n=100$, $n=230$, $n=250$), small AUROC gaps are
descriptive unless they recur across signal families or match the qualitative
traces.
No multiple-comparison correction is applied because the heatmap is not used
for family-wise significance testing.
The two primary headline contrasts instead carry paired-bootstrap intervals;
the largest, the adaptive-versus-strict SD-UQ collapse of $+0.52$ AUROC at
$n=196$, is far from zero, whereas gaps near $0.05$ are not interpreted.
Appendix Tables~\ref{tab:run_artifacts}, \ref{tab:effective_denominators},
\ref{tab:hotpot_fullwiki_stress}, and \ref{tab:headline_ci_output} document
run configurations, answered counts, GPS usable counts, the HotpotQA FullWiki
snapshot, and bootstrap intervals.
Retrieval-state AUROC with fewer than roughly fifty usable rows is treated as
trace-level evidence.
GPS AUROC is always conditional on rows for which GPS is defined; linking
failures and other abstentions are reported through usable/answered
denominators and audit-rule coverage, not hidden inside the AUROC.

The headline tables report family representatives: DSE, SRE-UQ, VN-Entropy,
and SD-UQ for answer state; SEU for evidence state; GPS for retrieval state.
The P(True)-style proxy is omitted because it is monotone in DSE here, and
SelfCheckGPT is retained in the artefacts.
GPS uses surface plus embedding-based soft answer-entity linking; its two
hyperparameters were calibrated on the RealMedQA development run and frozen
elsewhere.
All GPS values were recomputed by post-hoc replay on saved answer logs without
rerunning generation, retrieval, or entity linking
(Appendix~\ref{sec:appendix_reproducibility}).
PubMedQA-style yes/no answers remain outside GPS's intended scope because they
expose no linkable answer entities.

\subsection{Retrieval and KG Configuration}

Both systems use \texttt{all-MiniLM-L6-v2}.
Vanilla RAG performs direct vector retrieval over chunks.
KG-RAG performs entity-first retrieval, graph expansion, and context assembly
from graph-linked chunks, paths, and entities, with dense retrieval and
retriever-first graph expansion as fallbacks when entity anchoring is weak.
Each KG-RAG response records its route
(\texttt{entity\_first}, \texttt{rfge}, or \texttt{semantic\_only}) and route
reason.

Dataset-scoped KGs are built passage-wise, preserving question and passage
provenance and avoiding artificial cross-question chunking.
Each artefact retains chunk identifiers, passage titles where available,
relation paths, route labels, and relation-anchor text when supplied by the KG
extractor.
Dense retrieval can log passages and scores; the KG layer additionally exposes
matched entities, typed triples, graph paths, relation anchors, and abstention
events.
The traceability claim is therefore structured provenance, not provenance in
the abstract.
Traversal depth is dataset-specific: $h=2$ for PubMedQA, RealMedQA,
HotpotQA, HotpotQA FullWiki, and 2Wiki; $h=4$ for MuSiQue.
During iterative decomposition, each sub-question uses a bounded local graph
expansion cap so that overall chain length and per-step search breadth are not
confounded.
